\section{Facilities and Equipment Description}

The proposed research will be conducted within the California State University, Long Beach (CSULB) Hung Family College of Engineering and the Department of Mechanical and Aerospace Engineering (MAE). The College of Engineering emphasizes project-based learning, hands-on laboratory practice, and strong industry partnerships, providing an applied engineering environment well suited to the proposed hypersonic guidance, modeling, simulation, and validation effort. The MAE Department offers undergraduate and graduate programs in Aerospace Engineering and Mechanical Engineering and supports student training through a combination of theoretical instruction, laboratory practice, and project-based engineering activities.

\subsection{College and Department Environment}

CSULB MAE provides an appropriate academic setting for the proposed work in hypersonic guidance, navigation, and control. The proposed project will primarily require computational modeling, trajectory simulation, algorithm development, reinforcement learning training, uncertainty analysis, and technical reporting. These activities can be supported by PI-managed computing resources, university computing resources, department-supported software environments, and student research workstations.

The College of Engineering and MAE Department also provide access to complementary instructional and research laboratories that support aerospace systems, dynamics and control, robotics, manufacturing, materials, and experimental engineering. While the proposed Phase I effort is expected to emphasize simulation-based research rather than hardware procurement or flight testing, these facilities provide a strong environment for student training, future prototype development, and follow-on validation activities.

\subsection{Relevant MAE Laboratory Capabilities}

Relevant MAE laboratory capabilities include:

\begin{itemize}
  \item \textbf{Dynamics, Control, and Robotics Laboratory:} Supports instructional and research activities in dynamic system modeling, control system design and implementation, robotics modeling, and robotics control. Relevant equipment includes control plants, robotic manipulators, indoor 3D positioning systems, and depth cameras. This laboratory is relevant to the proposed project through its emphasis on dynamic systems, feedback control, sensing, and experimental validation of control concepts.
  \item \textbf{DENSO Design and Manufacturing Laboratory:} Provides facilities for computer-aided design, manufacturing, non-destructive testing, CNC machining, 3D scanning, and additive manufacturing. Although the proposed effort is primarily computational, this lab may support future hardware-oriented demonstrations, rapid prototyping, or educational activities connected to guidance system concepts.
  \item \textbf{Multiphase and Complex Fluids Flow Laboratory:} Supports computational and physics-based modeling of complex fluid mechanics problems. This capability is relevant to future extensions involving atmospheric modeling, aerodynamic uncertainty characterization, or higher-fidelity simulation environments.
  \item \textbf{Polymer and Advanced Composite Laboratory:} Supports development and characterization of lightweight materials and processes for aerospace and related industries. While not central to the proposed guidance algorithm development, this facility provides broader aerospace engineering context and potential support for future vehicle or test article studies.
  \item \textbf{Impact Group Engineering Research Laboratory:} Supports experimental and numerical studies of dynamic material and structural behavior under extreme environments such as shock, impact, and blast loading. This capability may be relevant to future research involving high-speed vehicle survivability, structural response, or validation under demanding environments.
\end{itemize}

\subsection{Computational and Software Resources}

The proposed research will rely on computational tools for trajectory propagation, numerical predictor-corrector guidance, multi-model simulation, reinforcement learning training, statistical analysis, and visualization. The project team expects to use appropriate engineering and scientific computing environments such as MATLAB/Simulink, Python, Julia, and C/C++, as needed. Version control will be used for source code, simulation configuration files, documentation, and analysis scripts.

Computational resources will be used to generate dispersed atmospheric and aerodynamic cases, evaluate baseline NPCG performance, train RL policies, compare guidance architectures, and quantify robustness against model uncertainty and dynamic PIP threat scenarios. If additional computing capacity is required for reinforcement learning training or large simulation campaigns, the team will use university-supported computing resources, PI-managed workstations, or other allowable computing resources identified during project execution.

\subsection{Suitability for the Proposed Effort}

The available CSULB College of Engineering and MAE environment is suitable for the proposed Phase I research because the effort is centered on analytical modeling, guidance algorithm development, simulation-based validation, and student-supported research. The combination of aerospace engineering programs, dynamics and control capability, computational modeling resources, and hands-on engineering laboratories provides a practical foundation for the proposed work and for future transition-oriented validation activities.
